\documentclass[../../../include/open-logic-section]{subfiles}

\begin{document}

\olfileid{sth}{replacement}{strength}
\olsection{प्रतिस्थापनाचे सामर्थ्य}

प्रतिस्थापनाच्या सामर्थ्याविषयी एक सोपे निरीक्षण आधी पाहू: $\Z$ च्या पुढे गेल्याशिवाय $\omega + \omega$ इतकी किंवा त्याहून मोठी कोणतीही फोन न्यूमन क्रमसूचक संख्या अस्तित्वात आहे, हे सिद्ध करता येत नाही.

का, याचा आराखडा असा. $\ZF$ मध्ये काम करताना $V_{\omega+\omega}$ हा संच पाहू. हा संच $\Z$ च्या एका \emph{प्रतिमानाचा} प्रांत म्हणून काम करतो. हे समजण्यासाठी सूत्राच्या \emph{सापेक्षीकरणाचा} संकेत ठरवू:
\begin{defn}\ollabel{formularelativization}
	कोणताही संच $M$ आणि कोणतेही सूत्र $\phi$ दिल्यास, $\phi$ मधील सर्व संख्यापक $M$ पुरते मर्यादित केल्यावर मिळणारे सूत्र $\phi^M$ असे मानू. म्हणजे ``$\lexists[x]$'' ऐवजी ``$(\lexists[x \in M])$'' आणि ``$\lforall[x]$'' ऐवजी ``$(\lforall[x \in M])$'' लिहू.
\end{defn}\noindent
$\Z$ च्या प्रत्येक स्वयंसिद्धक $\phi$ साठी $\ZF \vdash
\phi^{V_{\omega+\omega}}$ हे दाखवता येते. पण \olref[spine][rank]{ordsetrankalpha} वरून $\omega+\omega$ ही संख्या $V_{\omega+\omega}$ मध्ये \emph{नसते}. म्हणून $\omega+\omega$ अस्तित्वात नसण्याशी $\Z$ सुसंगत आहे.

म्हणूनच \olref[ordinals][replacement]{sec} मध्ये आपण म्हटले की \olref[ordinals][ordtype]{thmOrdinalRepresentation} प्रतिस्थापनाशिवाय सिद्ध करता येत नाही. कारण $\Z$ मध्येच असा स्पष्ट सुक्रमण-संबंध परिभाषित करणे सोपे आहे की ज्याचा क्रमप्रकार सहजपणे $\omega+\omega$ \emph{असायला हवा}. नैसर्गिक संख्यांवरील पुढील क्रम दाखवताना \olref[ordinals][idea]{sec} मध्ये आपण याचे अनौपचारिक उदाहरण दिले होते:
\begin{align*}
	n \lessdot m \text{ तेव्हाच आणि तेव्हाच }&\text{एकतर }n < m\text{ आणि }m-n\text{ सम आहे,}\\
	& \text{किंवा $n$ सम आणि $m$ विषम आहेत.}
\end{align*}
पण $\omega+\omega$ अस्तित्वात नसेल, तर हे सुक्रमण कोणत्याही क्रमसूचक संख्येशी समरूपी नसते. म्हणून $\Z$ वरून \olref[ordinals][ordtype]{thmOrdinalRepresentation} \emph{सिद्ध होत नाही}.

उलट पाहिल्यास: प्रतिस्थापनामुळे $\omega+\omega$ चे अस्तित्व, आणि म्हणून $V_{\omega+\omega}$ चे अस्तित्व, सिद्ध करता येते. एवढेच नव्हे. आपण परिभाषित करू शकणाऱ्या \emph{कोणत्याही} सुक्रमणासाठी \olref[ordinals][ordtype]{thmOrdinalRepresentation} सांगते की त्याच्याशी समरूपी अशी काही $\alpha$ असते; म्हणून $V_\alpha$ अस्तित्वात असते. अशा थेट अर्थाने प्रतिस्थापनामुळे संचांचा पदानुक्रम \emph{फार उंच} असण्याची हमी मिळते.

पुढील काही विभागांत आणि नंतर \olref[card-arithmetic][fix]{sec} मध्ये प्रतिस्थापनामुळे पदानुक्रम नेमका \emph{किती} उंच जातो याची अधिक चांगली कल्पना येईल. सध्या एवढाच साधा मुद्दा: प्रतिस्थापनाचे समर्थन \emph{खरोखर आवश्यक} आहे!

\end{document}
